\section{Nombre: Fantasma morado}   \label{per:fantasmaM}
\subsection{Descripción:}
Este enemigo nace de las almas que no quedaron varadas en los niveles del Mictlán al no poder terminar los desafíos. Tiene la forma de un cráneo, en los orificios de los ojos del cráneo se puede ver un brillo espectral. El cráneo está decorado con diferentes patrones de color morado y verde en la zona superior del cráneo, la frente, los ojos y la boca. De la parte trasera del cráneo sale fuego cuya paleta de colores es un degradado de morado oscuro y guinda, siendo el morado oscuro el predominante.  
\subsection{Status:}
\begin{itemize}
	\item Enemigo normal.
\end{itemize}
\subsection{Imagen}
Ver figura \ref{fig:fantasmaM}.
\begin{figure}
	\centering
	\includegraphics[height=0.2 \textheight]{Imagenes/fantasmaMorado}
	\caption{Concepto de diseño del Fantasma morado.}
	\label{fig:fantasmaM}
\end{figure}
\subsection{Encuentro:}
En niveles 2, 4, 5, 7, 9.
\subsection{Habilidades:}
\begin{itemize}
	\item Embestida (ver apartado \ref{hab.embestida}).
	
\end{itemize}
\subsection{Patrón de ataque:}
El enemigo realizara un recorrido vertical del punto A al B y del B al A. Este ciclo se repetirá tres veces, después se ejecutará la habilidad de embestida (ver apartado \ref{hab.embestida}). Todo este comportamiento se repetirá de manera cíclica.
\subsection{Bloques de animación:}
	\begin{itemize}
		\item Animación fuego.
	\end{itemize}